\section{Numerical evaluation}\label{app:numerical-evaluation}\label{sec:numerics}

The posterior construction requires three numerical operations. We
first prepare the ordinary Gaussian posterior, then evaluate the
shape threshold whenever the sampler proposes a coefficient vector,
and finally draw the omitted coefficient conditionally on each
retained vector. The Gaussian preparation is performed once, whereas
threshold evaluation is repeated throughout sampling. The main
numerical issue is therefore to bound a maximum over the entire
constraint interval reliably and to use the same bound in the
marginal density and conditional reconstruction.

\subsection{Gaussian preparation and the common algorithm}
\label{subsec:gaussian-preparation}

The preparation separates the coefficients that will be sampled
from the scalar that will be integrated out. Let $\boldsymbol A$ be
the invertible linear map from the original coefficients to
$(\boldsymbol g,\gamma)$ under the applicable conditions in
\cref{cor:positive,cor:monotone,cor:convex}, or the common map in
\cref{cor:combinations} for a joint restriction. The ordinary posterior
then has mean $\boldsymbol A\boldsymbol m$ and covariance
$\boldsymbol A\boldsymbol V\boldsymbol A^T$. Denote the mean blocks
by $\boldsymbol m_g$ and $m_\gamma$, the marginal covariance of
$\boldsymbol g$ by $\boldsymbol V_g$, the variance of $\gamma$ by
$v_\gamma$, and their cross-covariance by
$\boldsymbol v=\Cov(\boldsymbol g,\gamma\mid\boldsymbol y)$.
Gaussian conditioning gives
\begin{equation}\label{eq:gaussian-blocks}
 \mu(\boldsymbol g)=m_\gamma+
      \boldsymbol v^T\boldsymbol V_g^{-1}(\boldsymbol g-\boldsymbol m_g),
 \qquad s^2=v_\gamma-\boldsymbol v^T\boldsymbol V_g^{-1}\boldsymbol v.
\end{equation}
These formulas give the conditional mean and variance of $\gamma$
for any value of $\boldsymbol g$. To put the sampled coefficients
on a common Gaussian scale, choose an invertible factor
$\boldsymbol B\boldsymbol B^T=\boldsymbol V_g$ and set
$\boldsymbol g(\boldsymbol z)=\boldsymbol m_g+\boldsymbol B\boldsymbol z$.
Then \cref{eq:gaussian-blocks} gives
$\mu(\boldsymbol g(\boldsymbol z))=m_\gamma+\boldsymbol\lambda^T\boldsymbol z$,
where $\boldsymbol\lambda=\boldsymbol B^{-1}\boldsymbol v$ is
computed by a linear solve. Under the ordinary posterior,
$\boldsymbol z$ has a standard normal distribution. The shape
restriction then enters through the survival weight in
\cref{eq:whitened-target}.

When computing $s^2$, subtracting two nearly equal terms in
\cref{eq:gaussian-blocks} can lose precision. The equivalent formula
$s^2=(\boldsymbol d^T\boldsymbol V^{-1}\boldsymbol d)^{-1}$,
where $\boldsymbol d$ is the column of $\boldsymbol A^{-1}$
corresponding to $\gamma$, avoids this subtraction.
The quadratic form is the $\gamma$ diagonal entry of the transformed
precision matrix and can be evaluated using the factored posterior
precision without forming $\boldsymbol V^{-1}$ explicitly.

The ordinary Gaussian posterior can itself be prepared using matrix
factorizations and linear solves. For a general noise covariance
$\boldsymbol\Sigma$, solves with its Cholesky factor transform the
observation errors to have identity covariance. We apply the same transformation
to the design and to the response centered at
$\boldsymbol X\boldsymbol m_0$, then scale the design columns by
the prior standard deviations. In these coordinates, the prior
precision is the identity. Stacking the transformed design above an
identity matrix and applying a QR factorization therefore gives a
triangular factor of the posterior precision. Linear solves supply
the posterior mean and covariance factors, after which the full
transformation $\boldsymbol A\boldsymbol V\boldsymbol A^T$
preserves the dependence introduced by the coefficient map.

\cref{alg:posterior} separates the repeated marginal-density
evaluation from reconstruction after sampling. Both stages use the
same threshold rule, tolerance, and precision policy, so they describe
the same conservative conditioning event. A gradient-based sampler
also needs a vector to guide its proposals, called the force below.
We obtain it from the envelope derivative in
\cref{subsec:threshold-derivatives}, which differentiates the
maximized objective without differentiating the numerical search
for its maximizing location.

\begin{algorithm}[htbp]
\caption{Sign-constrained marginal sampling and reconstruction.
For a joint event satisfying \cref{prop:joint}, use its common
coefficient map and replace $\tau_r$ by $\tau_A^\varepsilon$
throughout, including the upper threshold used in both stages.}
\label{alg:posterior}
\begin{algorithmic}[1]
\Require Gaussian model, event $\mathcal C_r$ or $\mathcal C_A^\varepsilon$, admissible map, interval $I$
\Require Absolute tolerance $\delta>0$ and precision budget
\State Form the ordinary posterior and apply the invertible shape-specific map
\State Prepare $\boldsymbol m_g,\boldsymbol B,m_\gamma,\boldsymbol\lambda,s$
\Function{Marginal evaluation}{$\boldsymbol z$}
 \State Set $\boldsymbol g=\boldsymbol m_g+\boldsymbol B\boldsymbol z$
 \State Set $\mu=m_\gamma+\boldsymbol\lambda^T\boldsymbol z$
 \State Bound $\tau_r(\boldsymbol g)$ between $\ell$ and $u$ with $u-\ell\leq\delta$, or report failure
 \State Set $\widetilde\tau_r=u$ and $\widetilde\alpha=(\widetilde\tau_r-\mu)/s$
 \State Return $-\|\boldsymbol z\|^2/2+\log\Phi^{\mathrm c}(\widetilde\alpha)$ and the envelope force
\EndFunction
\State Use the marginal sampler to retain draws of $\boldsymbol z$
\For{each retained $\boldsymbol z$}
 \State Set $\boldsymbol g=\boldsymbol m_g+\boldsymbol B\boldsymbol z$ and $\mu=m_\gamma+\boldsymbol\lambda^T\boldsymbol z$
 \State Evaluate the same upper-threshold rule and set $\widetilde\alpha=(\widetilde\tau_r-\mu)/s$
 \State Draw $U\sim\operatorname{Uniform}(0,1)$ independently
 \State Compute $\eta=-\Phi^{-1}(U\Phi^{\mathrm c}(\widetilde\alpha))$ using log probabilities
 \State Set $\gamma=\widetilde\tau_r+s(\eta-\widetilde\alpha)$ and invert the coefficient map
 \State Evaluate the reconstructed curve and requested functionals
\EndFor
\end{algorithmic}
\end{algorithm}

\subsection{Conservative threshold evaluation}
\label{subsec:threshold-computation}

The threshold is a maximum over the entire constraint interval.
Evaluating its objective at finitely many locations gives a lower
bound, while an upper bound must also control values between those
locations. The calculation therefore maintains both bounds and
refines them until their difference is within the requested tolerance.

\paragraph{Bounds used for conditioning.}
The routine takes the remaining coefficients in their declared order,
the shape case and basis geometry, a positive absolute tolerance
$\delta$ in coefficient units, and an arithmetic precision limit.
A successful evaluation returns
\[
 \ell\leq\tau_r(\boldsymbol g)\leq u,
 \qquad u-\ell\leq\delta,
\]
together with an approximate maximizing location for the envelope
force. Choosing $\widetilde\tau_r=u$ then ensures
$0\leq\widetilde\tau_r-\tau_r\leq\delta$. The resulting lower bound on
$\gamma$ can be slightly too large, but cannot admit a curve that
violates the constraint. If the requested enclosure cannot be
obtained, the routine reports failure.

\paragraph{Locating and evaluating extrema.}
For each individual constraint, the maximum occurs at an endpoint
or an interior stationary location.
\cref{cor:positive,cor:monotone,cor:convex} express the interior
candidates as roots of a polynomial in the cosine coordinate $w\in J$.
Real-root isolation finds disjoint intervals containing every
distinct real root and establishes that none have been omitted
\citep{johansson2017arb}. The represented coefficients are treated as
exact rationals, allowing repeated factors to be removed before
isolation while retaining roots of even multiplicity. This arithmetic
also detects reduced degree without discarding small nonzero
coefficients. A constant stationary polynomial requires only endpoint
evaluations, including when it is identically zero.

Interval arithmetic then encloses the objective at the endpoints
and over each root interval. Rounding is directed outward so that
these enclosures include arithmetic error. The largest candidate
upper bound gives $u$, while valid lower bounds come from endpoints
and roots known to belong to $J$. Refining the candidate enclosures
reduces the gap between these bounds until it is at most $\delta$.

Roots near the boundary of $J$ need additional care. A root interval
known to lie outside $J$ is discarded. One crossing the boundary
remains an upper-bound candidate, but supplies a lower bound only
after its root is shown to belong to $J$. The calculation can still
stop without resolving that location if the overall gap is already
within $\delta$. For quotient objectives, division also requires a
strictly positive interval lower bound for the denominator.

\paragraph{Evaluation for the individual constraints.}
For positivity, isolate the roots of $P_{0,\boldsymbol g}$ and
evaluate the objective $-f_{\boldsymbol g}$ over the corresponding
physical locations. Interval evaluation of this trigonometric
expression encloses the candidate values.

For monotonicity, the objective is $-P_{\boldsymbol g}/D$, with
stationary polynomial
\[
 \frac{\mathrm{d}P_{\boldsymbol g}}{\mathrm{d}w}D
 -P_{\boldsymbol g}\frac{\mathrm{d}D}{\mathrm{d}w}.
\]
The polynomials in \cref{eq:monotone-polynomials} use every spectral
coefficient, with $g_2=0$ and the intercept excluded, and satisfy
$D\geq1$ on $J$. The lidar implementation first bounds the equivalent
trigonometric expression, retaining all harmonics, and uses root
isolation as a fallback.

For convexity, the polynomial objective $Q_{\boldsymbol g}$ can be
bounded before finding its stationary points. Map $J$ affinely to
$[0,1]$ and express the polynomial in the Bernstein basis. These
basis functions are nonnegative and sum to one, so every polynomial
value is a weighted average of the coefficients and lies between
their smallest and largest values. Ball arithmetic, which represents
each quantity by a midpoint and an error radius, encloses the change
of basis and the cosine endpoints.

De Casteljau subdivision rewrites the polynomial on smaller
intervals, with outward rounding preserving the bounds. Repeatedly
subdividing the interval with the largest upper bound refines $u$,
while endpoint and midpoint evaluations improve $\ell$.
If the gap still exceeds $\delta$ when the subdivision budget is
reached, isolate the roots of the exact rational derivative of
$Q_{\boldsymbol g}$ and increase arithmetic precision as needed.

\paragraph{Joint restrictions.}
For a signed joint event, use a common coefficient map satisfying
\cref{prop:joint}. Every selected order then shares the same direction
$h$ and remaining coefficients $\boldsymbol g$, so its quotient
objective must be formed from \cref{eq:joint-threshold} under this map.
Subdivide $I$ and bound each quotient on every subinterval, verifying
a positive lower bound for its signed denominator before division.
Evaluations at points in $I$ supply lower bounds.

If $[\ell_r,u_r]$ encloses the maximum for order $r\in A$, then
\[
 \max_{r\in A}\ell_r
 \leq\tau_A^\varepsilon(\boldsymbol g)
 \leq\max_{r\in A}u_r.
\]
The width of this joint enclosure is at most
$\max_{r\in A}(u_r-\ell_r)$, so requiring each width to be at most
$\delta$ suffices. The reductions in
\cref{cor:combinations,tab:constraint-combinations} replace some
continuum searches by endpoint evaluations. Enclose each retained
maximum and endpoint quotient with outward rounding, then take the
maxima of their lower and upper bounds. The same width guarantee
applies. The positive--convex and negative--concave pairs retain the
continuum value threshold $M_0$ alongside the curvature threshold
$M_2$.

\paragraph{Precision and error control.}
These guarantees concern the numerical coefficients and geometry
supplied to the routine, treated as exact inputs. They exclude errors
from the preceding Gaussian preparation. Applying \cref{prop:error}
requires the threshold-error bound to hold almost everywhere under
the ordinary Gaussian law.

Fixing the precision policy and the rules for choosing among
candidate locations makes repeated evaluation deterministic.
Nonfinite inputs and enclosures unresolved at the precision limit
are reported as failures. The Gaussian survival weight is positive
at every finite coefficient vector, so assigning zero density to an
unresolved evaluation would impose an additional restriction.
Such failures therefore cannot be silently treated as rejected
coefficient values.

\subsection{Threshold derivatives}
\label{subsec:threshold-derivatives}

Gradient-based sampling requires the derivative of the log marginal
density with respect to the sampled coefficients. The shape restriction
enters through the Gaussian survival weight in
\cref{eq:whitened-target}, so this calculation depends on how the
threshold changes with the coefficients. The maximum representation
of the threshold gives these derivatives from its maximizing locations,
without differentiating the numerical search used to find them.

\paragraph{Slopes at maximizing locations.}
At a fixed $x$, the threshold objective is linear in
$\boldsymbol g$. Under \cref{thm:finite}, write
\[
 \tau_r(\boldsymbol g)
 =\max_{x\in I}\boldsymbol\beta_r(x)^T\boldsymbol g,
\]
where $\boldsymbol\beta_r(x)$ contains the coefficient multipliers,
ordered as the entries of $\boldsymbol g$. The positive denominator
and compact interval make this vector continuous and bounded in $x$.
Consequently, $\tau_r$ is convex, positively homogeneous, and
Lipschitz, with Lipschitz constant at most
$\max_{x\in I}\|\boldsymbol\beta_r(x)\|$.

Let $\mathcal A_r(\boldsymbol g)$ denote the locations attaining
the maximum, which we call active locations.
For a perturbation in a direction $\boldsymbol b\in\R^m$,
Danskin's theorem \citep[Exercise 6.7]{bertsekas2010convexchapter}
gives
\[
 \left.\frac{\mathrm{d}}{\mathrm{d}u}
       \tau_r(\boldsymbol g+u\boldsymbol b)\right|_{u=0+}
 =\max_{x\in\mathcal A_r(\boldsymbol g)}
       \boldsymbol\beta_r(x)^T\boldsymbol b.
\]
The theorem applies because $I$ is compact and the coefficient
gradients are continuous. Its formula compares the slopes at the
currently maximizing locations and takes the largest slope in the
chosen direction.

The subdifferential consists of all convex combinations of the active
vectors $\boldsymbol\beta_r(x)$. A gradient exists exactly when all
these vectors agree, in which case their common value is the column
gradient. Thus a unique maximizing location suffices for
differentiability. Several locations also permit a gradient if they
give the same vector, while different active vectors produce a kink.

\paragraph{Individual and joint constraints.}
For the individual constraints, the coefficient vectors follow
directly from the objectives in
\cref{cor:positive,cor:monotone,cor:convex}. With
$w(x)=\cos(\omega_1(x-c+L))$, their spectral components are
\[
 \begin{aligned}
  \beta_{0,j}(x)&=-\phi_{j,L}(x),
       &&1\leq j\leq m,\\
  \beta_{1,j}(x)&=-jT_j(w(x))/D(w(x)),
       &&1\leq j\leq m,\quad j\ne2,\\
  \beta_{2,j}(x)&=j^2U_{j-1}(w(x)),
       &&2\leq j\leq m.
 \end{aligned}
\]
The intercept components in the monotonicity and convexity coordinates
are zero. Where the threshold is differentiable, these formulas
evaluate its gradient at any maximizing location.

For the joint threshold in \cref{prop:joint}, form the vectors from
the common coefficient map so that
$\eta_r(x;\boldsymbol g)=\boldsymbol\beta_r(x)^T\boldsymbol g$.
An active candidate is now a pair $(r,x)$ attaining the joint
threshold. The directional derivative is the largest
$\boldsymbol\beta_r(x)^T\boldsymbol b$ among these pairs, and a
gradient exists exactly when their vectors agree. This includes ties
between different derivative orders. The reductions in
\cref{cor:combinations} preserve the threshold, so the same calculation
can use the retained continuum objectives and endpoint quotients.

\paragraph{Gradient of the marginal log density.}
Let $q_r(\boldsymbol z)$ be the exact conditioned density in whitened
coordinates from \cref{eq:whitened-target}. Define the standardized
threshold by
\[
 \alpha_r(\boldsymbol z)
 =\frac{\tau_r(\boldsymbol g(\boldsymbol z))
          -\mu(\boldsymbol g(\boldsymbol z))}{s},
\]
and let $\varphi_1$ denote the standard normal density.
At a point where the active coefficient vectors agree, choose any
maximizing location $x_\star$. The map
$\boldsymbol g(\boldsymbol z)=\boldsymbol m_g+\boldsymbol B\boldsymbol z$
transforms the threshold gradient to
$\boldsymbol B^T\boldsymbol\beta_r(x_\star)$, while the conditional
mean has gradient $\boldsymbol\lambda$. The chain rule therefore gives
the column gradients
\begin{equation}\label{eq:target-gradient}
 \begin{aligned}
 \frac{\partial\alpha_r}{\partial\boldsymbol z}
   &=\frac{\boldsymbol B^T\boldsymbol\beta_r(x_\star)-\boldsymbol\lambda}{s},\\
 \frac{\partial\log q_r}{\partial\boldsymbol z}
   &=-\boldsymbol z-
       \frac{\varphi_1(\alpha_r)}{\Phi^{\mathrm c}(\alpha_r)}
       \frac{\partial\alpha_r}{\partial\boldsymbol z}.
 \end{aligned}
\end{equation}
The term $-\boldsymbol z$ comes from the standard normal density,
while the second term differentiates the log survival probability.
Its weight $\varphi_1(\alpha_r)/\Phi^{\mathrm c}(\alpha_r)$ is the
inverse Mills ratio. The gradient of $\alpha_r$ accounts for changes
in both the threshold and the conditional mean of $\gamma$.

\paragraph{The numerical force.}
The numerical adapter selects an approximate maximizing location
deterministically, evaluates $\boldsymbol\beta_r$ there, and uses
$\widetilde\alpha=(\widetilde\tau_r-\mu)/s$ in the normal-tail factor.
Substituting these quantities in the envelope formula gives the force
used by the sampler.

The threshold tolerance controls objective values, but locations with
almost equal values can have substantially different coefficient
slopes. A narrow enclosure of the maximum therefore does not guarantee
an accurate force near a change of maximizing location. The numerical
log density uses the upper threshold, so the returned envelope force need
not be its exact gradient. \cref{app:stan-implementation} explains how
this force is used in sampling.

\subsection{Stable tail evaluation and reconstruction}
\label{subsec:tail-computation}

When the threshold lies far above the scalar conditional mean, its
Gaussian survival probability can be extremely small. Computing
$1-\Phi(\alpha)$ can then lose precision or return zero, even
though the probability is positive. We instead evaluate
$\log\Phi^{\mathrm c}(\alpha)$ directly as $\log\Phi(-\alpha)$ using a
normal log-tail routine. The inverse Mills ratio
in \cref{eq:target-gradient} is then
$\exp(\log\varphi_1(\alpha)-\log\Phi^{\mathrm c}(\alpha))$.

For reconstruction, draw $U\sim\operatorname{Uniform}(0,1)$
independently and use a normal quantile routine that accepts log
probabilities to solve
\[
 \log\Phi(-\eta)=\log U+\log\Phi^{\mathrm c}(\widetilde\alpha).
\]
In exact arithmetic this gives $\eta>\widetilde\alpha$ for
$0<U<1$, and
$\gamma=\widetilde\tau_r+s(\eta-\widetilde\alpha)$ equals the
truncated-normal draw in \cref{eq:reconstruction} with the upper
threshold substituted. Expressing the draw relative to the threshold
is useful when the conditional mean and a large positive normal
quantile would nearly cancel. A nonfinite quantile or a
computed $\eta<\widetilde\alpha$ is a reconstruction failure.
These arithmetic safeguards do not extend \cref{prop:error} to
rounding in Gaussian preparation or tail evaluation, since that
bound assumes exact Gaussian and tail operations.
